
\subsection{}
\label{XIII.4.7} Sea
$$
\xymatrix@C=1.2cm{
X \ar[d]_-f&\ar[l] X_U \ar[d]_-{f_U}& \ar[l]X_{\bar{s}}\ar[d] \\
S&\ar[l] U &\ar[l]\bar{s}
}
$$
un
\marginpar{423}
diagrama cuyos cuadrados son cartesianos, donde $S$ es un esquema conexo
por arcos (SGA~4 IX~2.12), $U$ un abierto conexo de~$S$ y~$\bar{s}$ un
punto geométrico de~$U$. Sean~$a$ un punto geométrico de~$X_{\bar{s}}$
y~$\LL$ un conjunto de números primos. Sea~$g$ una sección de~$f$ y
supongamos que se satisfacen las condiciones siguientes:
\begin{enumerate}
\item[a)] El morfismo $f$ es $0$-acíclico, localmente $0$-acíclico,
y, para todo recubrimiento étale $S'$ de~$S$ y todo recubrimiento
étale $E$ de~$X\times_S S'$, cociente de un recubrimiento galoisiano
cuyo grupo es un~$\LL$-grupo, $(E,f_{(S')})$ es cohomológicamente propio
relativamente a~$S'$ en dimensión $\leq 0$.
\item[b)] El morfismo $f_U$ es localmente $1$-asférico para~$\LL$ y,
para todo haz constante finito de~$\LL$-grupos~$F$ sobre~$X_U$, $(F,f_U)$
es cohomológicamente propio en dimensión $\leq 1$ y las fibras de
$\R^1 f_{U*}F$ son finitas.
\end{enumerate}

Se deduce entonces de~\Ref{XIII.4.1} y~\Ref{XIII.4.3} el diagrama
conmutativo siguiente, cuyas filas son exactas:
\begin{equation*}
\label{eq:XIII.4.7.0} \tag{\thesubsection.0} {}
\begin{array}{c}
\xymatrix{
1 \ar[r]& \pi_1^{\LL}(X_{\bar{s}},a) \ar[r]\ar@{=}[d]& \pi'_1(X_U,a) \ar[r]\ar[d]& \pi_1(U,a)\ar[r]\ar[d]& 1 \\
&\pi_1^{\LL}(X_{\bar{s}},a) \ar[r]&\pi'_1(X,a) \ar[r]& \pi_1(S,a) \ar[r]& 1
}
\end{array}
\end{equation*}

Gracias a la sección~$g$, se tienen morfismos
\ifthenelse{\boolean{orig}}
{$\pi_1(U,a)\to\pi'_1(X_U,a)$ $\pi_1(S,a)\to\pi'_1(X,a)$;}
{$$\pi_1(U,a)\to\pi'_1(X_U,a),\quad \pi_1(S,a)\to\pi'_1(X,a)\,\;$$}
de ellos se deduce un morfismo de la suma amalgamada de~$\pi_1(S,a)$ y
$\pi'_1(X_U,a)$ sobre~$\pi_1(U,a)$ en~$\pi'_1(X,a)$:
\begin{equation*}
\label{eq:XIII.4.7.1} \tag{\thesubsection.1} {}
\varphi\colon\pi=\pi_1(S,a)\coprod_{\pi_1(U,a)}\pi'_1(X_U,a)\to\pi'_1(X,a).
\end{equation*}
Supongamos que se satisface la condición siguiente:
\begin{enumerate}
\item[c)] Si $T=S-U$, se tiene $\prof \et_T(S)\geq 2$ (SGA~2 XIV~1.1).
\end{enumerate}
Entonces, el funtor que a un recubrimiento étale de~$S$ le hace
corresponder su restricción a~$U$ es plenamente fiel (SGA~2 XVI~1.4).
Se sigue que el morfismo $\pi_1(U,a)\to\pi_1(S,a)$ es sobreyectivo
(\Ref{V.6.9}), y del
\marginpar{424}
diagrama~\eqref{eq:XIII.4.7.0} se deduce que lo mismo ocurre con el
morfismo $\pi'_1(X_U,a)\to\pi'_1(X,a)$; con mayor razón, $\varphi$ es un
epimorfismo. Sea
\begin{equation*}
\label{eq:XIII.4.7.2} \tag{\thesubsection.2} {}
K=\Ker(\pi_1(U,a)\to\pi_1(S,a)).
\end{equation*}
El grupo~$\pi$ de~\eqref{eq:XIII.4.7.1} se identifica con el cociente de
\ifthenelse{\boolean{orig}}
{$\pi'_1(X_U,a)$ por el subgrupo normal cerrado engendrado por la imagen~$L$
de~$K$ en~$\pi'_1(X_U,a)$.}
{$\pi'_1(X_U,a)$ por el subgrupo normal cerrado~$L$ engendrado por la imagen
de~$K$ en~$\pi'_1(X_U,a)$.}
Consideremos $\pi'_1(X_U,a)$ como producto
semidirecto de~$\pi_1(U,a)$ por~$\pi_1^\LL (X_{\bar{s}},a)$. El grupo~$K$
actúa entonces por automorfismos interiores sobre
$\pi_1^\LL (X_{\bar{s}},a)$, y el cociente
$\pi=\pi'_1(X_U,a)/L$ se identifica con el producto semidirecto de
$\pi_1(U,a)/K=\pi_1(S,a)$ por el grupo
$\pi_1^\LL(X_{\bar{s}},a)_K$ de coinvariantes de
$\pi_1^{\LL}(X_{\bar{s}},a)$ bajo la acción de~$K$. Finalmente, se tiene un
epimorfismo
\begin{equation*}
\label{eq:XIII.4.7.3} \tag{\thesubsection.3} {}
\varphi\colon\pi=\pi_1^\LL(X_{\bar{s}},a)_K\cdot\pi_1(S,a)\to\pi'_1(X,a).
\end{equation*}
\label{indnot:ml}\oldindexnot{$\pi_1^\LL(X_{\bar{s}},a)_K$|hyperpage}%

La proposición siguiente da condiciones bajo las cuales el morfismo
$\varphi$ es un isomorfismo.
